\documentclass[../../../../main.tex]{subfiles}
\begin{document}

\begin{flushright}
\footnotesize\textit{【自動文字起こし・要確認】}
\end{flushright}

[解] $f(x)$ を $f$ などと略記する

\begin{flushright}
$\left[ \text{よく考えると、わざわざ (2) で } (A,B)=(1,3) \text{ を排除する必要なかったね…} \right]$
\end{flushright}

(1) $g(x), h(x)$ は整数係数だから、$g(0), h(0) \in \mathbb{Z} \cdots \text{①}$.
又、$f(0) = g(0)h(0)$ から $1 = g(0)h(0) \cdots \text{②}$
①, ②から $(g(0), h(0)) = (1, 1), (-1, -1)$ となるが、いずれの場合も $g(0) = h(0)$ となる 終

(2) 題意の時、対称性から $(g\text{の次数}) \le (h\text{の次数})$ とすると、
$(A, B) = (1, 3), (2, 2)$ のいずれかである。
前者の時、$f$ の最高次係数が $1$ なことと (1) から
$g(x) = \pm x \pm 1$ となるが、$f(\pm 1) \neq 0$ と因数定理から、$f(x) = g(x)h(x)$ に矛盾。
よって $(A, B) = (2, 2)$ に限られる。以下 $f(m), g(m) \in \mathbb{Z}$ ($m \in \mathbb{Z}$) に注意する。

\[
\begin{cases}
g(1)h(1) = 1 \\
g(a)h(a) = 1 \\
g(b)h(b) = 1
\end{cases}
\]
より、(1) と同じく $g(1)=h(1), g(a)=h(a), g(b)=h(b)$ である。
$g, h$ は共に2次式であることと、異なる3つの値で $g, h$ の値が一致することから、$g(x) = h(x)$ となる 終

(3) $g=h$ の時
\[ (g+1)(g-1) = x(x-a)(x-1)(x-b) \cdots \text{③} \]
たとえば $(a, b) = (-1, -2)$、$g(x) = x^2+x-1$ とすれば
\[ (g+1)(g-1) = x(x+1)(x+2)(x-1) \]
となって③をみたすから
$(a, b) = (-1, -2)$ が1例である。終

\end{document}
